% Research candidate for P06.2.  Not authorized for live-book use before
% independent exact-byte audit and production Rethlas verification.

\begin{proof}[Proof of \cref{prop:p062-finite-flat-continuous-trace}]
The assertions are local on \(X\), so fix an affinoid domain \(U\subset X\).
Its inverse image \(Y_U\) is affinoid.  By the smooth
Stone--Weierstrass theorem of Chambert--Loir--Ducros
\cite[Theorem~15.2.6]{CLD25}, there are ordinary CLD-smooth functions
\(h_j\) on \(Y_U\) which converge uniformly to \(h|_{Y_U}\).
For every \(j\), put
\[
 H_j(x):=\int_{Y_x}h_j|_{Y_x},\qquad x\in U.
\]
By \cite[Theorem~16.3.3(d)]{CLD25}, fibre integration in relative
dimension zero makes \(H_j\) \(G\)-smooth, hence continuous, on \(U\).
Every fibre has
calibrated mass \(d\).  Consequently
\begin{equation}\label{eq:p062-trace-uniform-bound}
 \lVert H_j-\operatorname{Tr}_p(h)\rVert_{C^0(U)}
 \leq d\,\lVert h_j-h\rVert_{C^0(Y_U)}.
\end{equation}
Thus the trace in \eqref{eq:p062-finite-trace-definition} is continuous.

For a smooth \(h_j\), the projection formula for fibre integration
\cite[Corollary~16.5.5]{CLD25} gives
\begin{equation}\label{eq:p062-smooth-trace-current}
                         p_*[h_j]=[H_j]
\end{equation}
on \(U\).  Uniform convergence against every compactly supported smooth
top-degree test form lets \(j\to\infty\) in
\eqref{eq:p062-smooth-trace-current}.  Together with
\eqref{eq:p062-trace-uniform-bound}, this proves
\eqref{eq:p062-finite-trace-current} on \(U\), and hence on \(X\).

Finally, direct image by a compact morphism commutes with
\(\mathrm d'\) and \(\mathrm d''\)
\cite[Section~7.7.8]{CLD25}.  Restricting
\eqref{eq:p062-finite-trace-current} to \(V\) therefore gives
\[
 \mathrm d'\mathrm d''[\operatorname{Tr}_p(h)]
   =p_*\bigl(\mathrm d'\mathrm d''[h]\bigr)=0.
\]
This proves \eqref{eq:p062-finite-trace-hessian}.
\end{proof}

\begin{proof}[Proof of \cref{prop:p062-arbitrary-field-disc}]
We treat the open and closed discs simultaneously until the final density
step.  Let \(y\) be a rigid point of the disc under consideration, let
\(E=\mathscr H(y)\), and let \(P\in K[T]\) be the monic irreducible
polynomial of \(T(y)\).  Write
\[
 P(T)=T^d+\sum_{j=0}^{d-1}c_jT^j,
 \qquad d=[E:K],
 \qquad s=|T(y)|.
\]
The absolute value of \(K\) has a unique extension to every finite
extension.  Hence every algebraic conjugate of \(T(y)\), counted with its
inseparable multiplicity, has absolute value \(s\).  The elementary
symmetric functions of these conjugates give
\begin{equation}\label{eq:p062-minimal-polynomial-coefficient-bound}
                         |c_j|\leq s^{d-j}
                         \qquad(0\leq j<d).
\end{equation}

In the open case one has \(s<R\), and the monic polynomial \(P\) induces
a finite flat morphism of degree \(d\)
\begin{equation}\label{eq:p062-open-polynomial-map}
 p\colon D_R^-\longrightarrow D_{R^d}^- ,
 \qquad Z\longmapsto P(T).
\end{equation}
Indeed, \eqref{eq:p062-minimal-polynomial-coefficient-bound} shows that
\(|P(x)|<R^d\) if and only if \(|T(x)|<R\).  The morphism is the indicated
base change of the finite free polynomial morphism
\(\mathbf A_K^1\to\mathbf A_K^1\); freeness, with basis
\(1,T,\ldots,T^{d-1}\), also follows directly by division by the monic
polynomial \(P(T)-Z\).

In the closed case, \(s\leq R\), and the same argument gives a finite flat
degree-\(d\) morphism
\begin{equation}\label{eq:p062-closed-polynomial-map}
 p\colon D_R^+\longrightarrow D_{R^d}^+ .
\end{equation}
The fibre over the \(K\)-rational point \(0\) is
\(
 \mathcal M(K[T]/(P))=\{y\}
\)
and has calibrated degree \(d\).

Set \(H=\operatorname{Tr}_p(h)\).  By
\cref{prop:p062-finite-flat-continuous-trace}, \(H\) is continuous.  In the
open case it satisfies
\(
 \mathrm d'\mathrm d''[H]=0
\)
on \(D_{R^d}^-\).  In the closed case the same equation holds on the
absolute analytic interior of \(D_{R^d}^+\).  For the latter assertion,
the Gauss-fibre calculation below, applied with \(t=R\), shows that the
inverse image of the unique analytic boundary point \(\xi_{R^d}\) is the
unique analytic boundary point \(\eta_R\).  Thus the pullback of a test form
supported in the target interior is supported in the source interior, where
the current equation for \(h\) holds.

We next record the fibre at a suitable Gauss point.  In the open case choose
\begin{equation}\label{eq:p062-radius-above-rigid-point}
                              s<t<R,
\end{equation}
and put \(q=t^d\).  In the closed case make the same choice when \(s<R\),
and put \(t=R\), \(q=R^d\), when \(s=R\).  Let \(\eta_t\) be the Gauss
point of the source disc and \(\xi_q\) the Gauss point of the target disc.
Then
\begin{equation}\label{eq:p062-single-gauss-fibre}
                         p^{-1}(\xi_q)=\{\eta_t\},
 \qquad \deg_{\eta_t}(p)=d.
\end{equation}

For completeness, \eqref{eq:p062-single-gauss-fibre} follows directly from
Gauss valuations.  By
\eqref{eq:p062-minimal-polynomial-coefficient-bound}, the \(t\)-Gauss norm
of \(P\) is \(q=t^d\), and its initial form has degree \(d\).  If
\(U\) and \(V\) denote the homogeneous residues of \(T\) and \(Z\), the
initial equation of \(P(T)-Z\) is
\[
                      \operatorname{in}_t(P)(U)-V.
\]
The element \(\operatorname{in}_t(P)(U)\) is transcendental over the
graded residue field.  It follows first that composition with \(P\) carries
the \(t\)-Gauss valuation to the \(q\)-Gauss valuation, so
\(p(\eta_t)=\xi_q\).  Moreover, \(V\) is transcendental over the graded
residue field, and the displayed
polynomial is irreducible over its graded fraction field: equivalently,
the rational-function extension generated by \(U\) over that generated by
\(\operatorname{in}_t(P)(U)\) has degree \(d\).  This remains true in the
inseparable case.  The graded fundamental inequality therefore gives local
degree at least \(d\) at \(\eta_t\).  Since the finite flat morphism has total
degree \(d\), there is no other point in the fibre and the local degree is
exactly \(d\).  This proves \eqref{eq:p062-single-gauss-fibre} without
splitting \(P\) into geometric roots.

The zero fibre and \eqref{eq:p062-single-gauss-fibre} now give
\begin{equation}\label{eq:p062-trace-two-fibres}
                         H(0)=d\,h(y),
 \qquad                  H(\xi_q)=d\,h(\eta_t).
\end{equation}
Apply \cref{prop:p062-rational-centred-traces} to \(H\) on the target
disc.  Its radial identity, including the endpoint in the closed case,
and \eqref{eq:p062-trace-two-fibres} yield
\begin{equation}\label{eq:p062-rigid-to-gauss}
                              h(y)=h(\eta_t).
\end{equation}
Applying the same proposition to \(h\) on the source disc gives
\(h(\eta_t)=h(0)\).  Thus every rigid point \(y\) has the same value
\(h(0)\).

Rigid points are dense in every strict affinoid domain; see
\cite[Proposition~2.1.15]{Berk12}.  This applies directly to \(D_R^+\).
Applied locally on a strict-affinoid covering of \(D_R^-\), it also proves
density there.  Continuity of \(h\) proves constancy on the entire open or
closed disc.
\end{proof}
